\section{Exact Filter-Seven Localization}
\label{sec:exact-filter-seven-localization}

At the first nontrivial conditioned layer, exact residue order replaces the
quadratic capacity envelope by a constant boundary discrepancy, for every
finite integer interval $I$, after filters $2$, $3$, and $5$ have been
installed and immediately before filter $7$, over actual pre-filter-$7$
2-gap starts in an arbitrary integer interval. This is an arithmetic
localization theorem, not a density estimate.

The incoming starts occupy exactly three classes modulo $30$. Define
\begin{equation*}
F_7(x)=\mathbf1_{\{11,17,29\}\bmod30}(x),
\qquad
h_7(x)=\mathbf1_{\{0,5\}\bmod7}(x),
\end{equation*}

where $h_7$ marks the two endpoint-strike classes. The centered observable and
its interval sum are
\begin{equation*}
g_7(x)=F_7(x)\left(h_7(x)-\frac27\right),
\qquad
b_7(I)=\sum_{x\in I}g_7(x).
\end{equation*}

Because $30$ and $7$ are coprime, $g_7$ has period $210$. Its 21 admissible
start residues in increasing order are
\begin{equation*}
\begin{aligned}
&11,17,29,41,47,59,71,77,89,101,107,\\
&119,131,137,149,161,167,179,191,197,209.
\end{aligned}
\end{equation*}

Clearing the denominator gives weight $5$ at a harmful residue and $-2$ at a
harmless residue. In the order above the exact sequence is
\begin{equation*}
-2,-2,-2,-2,5,-2,-2,5,5,-2,-2,5,5,-2,-2,5,-2,-2,-2,-2,-2.
\end{equation*}

It contains six harmful and fifteen harmless terms, so
\begin{equation*}
6\cdot5+15\cdot(-2)=0
\qquad[\text{Complete-Period Cancellation}]
\end{equation*}

Starting from zero, its cumulative sums are
\begin{equation*}
\begin{aligned}
0,&-2,-4,-6,-8,-3,-5,-7,-2,3,1,-1,\\
&4,9,7,5,10,8,6,4,2,0.
\end{aligned}
\end{equation*}

Their minimum is $-8$ and maximum is $10$. Every non-wrapping consecutive
subsum is the difference of two cumulative sums. Every wrapping subsum is the
negative of its non-wrapping complement because the full sum is zero. Hence
every cyclic subsum $C$ satisfies
\begin{equation*}
\begin{aligned}
|C|
&\le10-(-8)
&&[\text{Cumulative-Sum Range}]\\
&=18.
&&[\text{Simplification}]
\end{aligned}
\end{equation*}

Now partition any integer interval $I$ into consecutive complete blocks of
length $210$ and one remainder of length less than $210$. The complete blocks
contribute zero, and the remainder selects one cyclic subsum. Therefore
\begin{equation*}
\begin{aligned}
|7b_7(I)|&\le18
&&[\text{Complete Blocks Cancel}]\\
|b_7(I)|&\le\frac{18}{7}.
&&\text{[Divide By }7\text{]}\quad\blacksquare\ \text{[Q.E.D.]}
\end{aligned}
\end{equation*}

The interval from residue $47$ through residue $161$ attains $18/7$, so the
constant is sharp. For the complete conditioned chain with
$r_0=5$ and $r_1=7$, write $P_m=A_{0,m}$. Then
\begin{equation*}
\begin{aligned}
w_1
&=A_{2,m}
&&[\text{By Definition}]\\
&=\frac{P_m}{a_0a_1}
&&[\text{Factor }A_{0,m}]\\
&=\frac{P_m}{(3/5)(5/7)}
&&[\text{Substitute }r_0=5,\ r_1=7]\\
&=\frac{7P_m}{3},
&&[\text{Simplification}]\\
\alpha_1
&=w_1\frac{r_1}{2(r_1-2)}
&&[\text{Energy Coefficient}]\\
&=\frac{49P_m}{30}.
&&[\text{Substitution}]
\end{aligned}
\end{equation*}

Therefore
\begin{equation*}
\begin{aligned}
\alpha_1b_7^2
&\le\frac{49P_m}{30}\left(\frac{18}{7}\right)^2
&&[\text{Sharp Interval Bound}]\\
&=\frac{54}{5}P_m.
&&[\text{Simplification}]
\end{aligned}
\end{equation*}

The saving comes from the exact ordered residue pattern. It controls one
fixed early layer and does not give a uniform bound for the growing family of
later coefficients.

{\raggedright
The derivation above proves the exact certificate and arbitrary-interval
bound. \par}
